\documentclass[10pt, letterpaper]{article}

% Packages:
\usepackage[
    ignoreheadfoot, % set margins without considering header and footer
    top=2 cm, % seperation between body and page edge from the top
    bottom=2 cm, % seperation between body and page edge from the bottom
    left=2 cm, % seperation between body and page edge from the left
    right=2 cm, % seperation between body and page edge from the right
    footskip=1.0 cm, % seperation between body and footer
    % showframe % for debugging 
]{geometry} % for adjusting page geometry
\usepackage[explicit]{titlesec} % for customizing section titles
\usepackage{tabularx} % for making tables with fixed width columns
\usepackage{array} % tabularx requires this
\usepackage[dvipsnames]{xcolor} % for coloring text
\definecolor{primaryColor}{RGB}{4, 110, 143} % define primary color
\definecolor{backgroundColor}{RGB}{255, 255, 240}
\definecolor{linkColor}{RGB}{77, 157, 224}
\definecolor{strongColor}{RGB}{21, 21, 21}
\definecolor{sourceColor}{RGB}{105, 114, 104}
\usepackage{enumitem} % for customizing lists
\usepackage{fontawesome5} % for using icons
\usepackage{amsmath} % for math
\usepackage[
    pdftitle={Giannis Papaioannou's CV},
    pdfauthor={Giannis Papaionnou},
    pdfcreator={LaTeX with RenderCV},
    colorlinks=true,
    urlcolor=linkColor
]{hyperref} % for links, metadata and bookmarks
\usepackage[pscoord]{eso-pic} % for floating text on the page
\usepackage{calc} % for calculating lengths
\usepackage{bookmark} % for bookmarks
\usepackage{lastpage} % for getting the total number of pages
\usepackage{changepage} % for one column entries (adjustwidth environment)
\usepackage{paracol} % for two and three column entries
\usepackage{ifthen} % for conditional statements
\usepackage{needspace} % for avoiding page brake right after the section title
\usepackage{iftex} % check if engine is pdflatex, xetex or luatex
\usepackage{pagecolor,lipsum}% http://ctan.org/pkg/{pagecolor,lipsum}
\newcommand\boldblack[1]{\textcolor{strongColor}{\textbf{#1}}}
\newcommand\textscc[1]{\textcolor{sourceColor}{\textsc{#1}}}
% Ensure that generate pdf is machine readable/ATS parsable:
\ifPDFTeX
    \input{glyphtounicode}
    \pdfgentounicode=1
    \usepackage[T1]{fontenc}
    \usepackage[utf8]{inputenc}
    \usepackage{lmodern}
\fi

\usepackage[default, type1]{sourcesanspro} 

% Some settings:
\AtBeginEnvironment{adjustwidth}{\partopsep0pt} % remove space before adjustwidth environment
\pagestyle{empty} % no header or footer
\setcounter{secnumdepth}{0} % no section numbering
\setlength{\parindent}{0pt} % no indentation
\setlength{\topskip}{0pt} % no top skip
\setlength{\columnsep}{0.15cm} % set column seperation
\makeatletter
\let\ps@customFooterStyle\ps@plain % Copy the plain style to customFooterStyle
\patchcmd{\ps@customFooterStyle}{\thepage}{
    \color{sourceColor}\textit{\small Page \thepage{} of \pageref*{LastPage}}
}{}{} % replace number by desired string
\makeatother
\pagestyle{customFooterStyle}

\titleformat{\section}{
    % avoid page braking right after the section title
    \needspace{4\baselineskip}
    % make the font size of the section title large and color it with the primary color
    \Large\color{primaryColor}
}{
}{
}{
    % print bold title, give 0.15 cm space and draw a line of 0.8 pt thickness
    % from the end of the title to the end of the body
    \textbf{#1}\hspace{0.15cm}\titlerule[0.8pt]\hspace{-0.1cm}
}[] % section title formatting

\titlespacing{\section}{
    % left space:
    -1pt
}{
    % top space:
    0.3 cm
}{
    % bottom space:
    0.2 cm
} % section title spacing

% \renewcommand\labelitemi{$\vcenter{\hbox{\small$\bullet$}}$} % custom bullet points
\newenvironment{highlights}{
    \begin{itemize}[
        topsep=0.10 cm,
        parsep=0.10 cm,
        partopsep=0pt,
        itemsep=0pt,
        leftmargin=0.4 cm + 10pt
    ]
}{
    \end{itemize}
} % new environment for highlights

\newenvironment{highlightsforbulletentries}{
    \begin{itemize}[
        topsep=0.10 cm,
        parsep=0.10 cm,
        partopsep=0pt,
        itemsep=0pt,
        leftmargin=10pt
    ]
}{
    \end{itemize}
} % new environment for highlights for bullet entries


\newenvironment{onecolentry}{
    \begin{adjustwidth}{
        0.2 cm + 0.00001 cm
    }{
        0.2 cm + 0.00001 cm
    }
}{
    \end{adjustwidth}
} % new environment for one column entries

\newenvironment{twocolentry}[2][]{
    \onecolentry
    \def\secondColumn{#2}
    \setcolumnwidth{\fill, 3.5 cm}
    \begin{paracol}{2}
}{
    \switchcolumn \raggedleft \secondColumn
    \end{paracol}
    \endonecolentry
} % new environment for two column entries

\newenvironment{flextwocolentry}[2][]{
    \onecolentry
    \def\secondColumn{#2}
    \setcolumnwidth{\fill, \fill}
    \begin{paracol}{2}
}{
    \switchcolumn \raggedleft \secondColumn
    \end{paracol}
    \endonecolentry
} % new environment for two column entries

\newenvironment{certtwocolentry}[2][]{
    \onecolentry
    \def\secondColumn{#2}
    \setcolumnwidth{\fill, 5.5 cm}
    \begin{paracol}{2}
}{
    \switchcolumn \raggedleft \secondColumn
    \end{paracol}
    \endonecolentry
} % new environment for two column entries
\newenvironment{threecolentry}[3][]{
    \onecolentry
    \def\thirdColumn{#3}
    \setcolumnwidth{1 cm, \fill, 4.5 cm}
    \begin{paracol}{3}
    {\raggedright #2} \switchcolumn
}{
    \switchcolumn \raggedleft \thirdColumn
    \end{paracol}
    \endonecolentry
} % new environment for three column entries

\newenvironment{flexthreecolentry}[3][]{
    \onecolentry
    \def\thirdColumn{#3}
    \setcolumnwidth{1 cm, \fill, \fill}
    \begin{paracol}{3}
    {\raggedright #2} \switchcolumn
}{
    \switchcolumn \raggedleft \thirdColumn
    \end{paracol}
    \endonecolentry
} % new environment for three column entries

\newenvironment{header}{
    \setlength{\topsep}{0pt}\par\kern\topsep\centering\color{primaryColor}\linespread{1.5}
}{
    \par\kern\topsep
} % new environment for the header

\newcommand{\placelastupdatedtext}{% \placetextbox{<horizontal pos>}{<vertical pos>}{<stuff>}
  \AddToShipoutPictureFG*{% Add <stuff> to current page foreground
    \put(
        \LenToUnit{\paperwidth-2 cm-0.2 cm+0.05cm},
        \LenToUnit{\paperheight-1.0 cm}
    ){\vtop{{\null}\makebox[0pt][c]{
        \small\color{sourceColor}\textit{Last updated in February 2026}\hspace{\widthof{Last updated in February 2026}}
    }}}%
  }%
}%

% save the original href command in a new command:
\let\hrefWithoutArrow\href

% new command for external links:
\renewcommand{\href}[2]{\hrefWithoutArrow{#1}{\ifthenelse{\equal{#2}{}}{ }{#2 }\raisebox{.15ex}{\footnotesize \faExternalLink*}}}


\begin{document}
\pagecolor{backgroundColor}
    \newcommand{\AND}{\unskip
        \cleaders\copy\ANDbox\hskip\wd\ANDbox
        \ignorespaces
    }
    \newsavebox\ANDbox
    \sbox\ANDbox{}

    \placelastupdatedtext
    \begin{header}
        \fontsize{30 pt}{30 pt}
        \textbf{Giannis Papaioannou}
        \boldblack{}

        \vspace{0.15 cm}
        \fontsize{20 pt}{20 pt}
        \textbf{Senior Python Developer}

        \normalsize
        \vspace{0.15 cm}
        \mbox{\color{linkColor}{\footnotesize\faMapMarker*}\hspace*{0.13cm}Athens, Greece}%
        \kern 0.25 cm%
        \AND%
        \kern 0.25 cm%
        \mbox{\hrefWithoutArrow{mailto:papaioannou.giannis@protonmail.com}{{\footnotesize\faEnvelope[regular]}\hspace*{0.13cm}papaioannou.giannis@protonmail.com}}%
        \kern 0.25 cm%
        \AND%
        \kern 0.25 cm%
        \mbox{\hrefWithoutArrow{tel:+30-6947884234}{{\footnotesize\faPhone*}\hspace*{0.13cm}+306947884234}}%
        \newline
        \kern 0.25 cm%
        \AND%
        \kern 0.25 cm%
        \mbox{\hrefWithoutArrow{https://linkedin.com/in/giannis-papaioannou/}{{\footnotesize\faLinkedinIn}\hspace*{0.13cm}giannis-papaioannou}}%
        \kern 0.25 cm%
        \AND%
        \kern 0.25 cm%
        \mbox{\hrefWithoutArrow{https://github.com/giannis-papaioannou/}{{\footnotesize\faGithub}\hspace*{0.13cm}giannis-papaioannou}}%
    \end{header}

    \vspace{0.3 cm - 0.3 cm}


    \section{Summary}
        \begin{onecolentry}
            Senior Software Engineer specializing in \texttt{Python}, with expertise in building scalable systems and solving complex technical challenges. Engineered robust solutions for data pipelines and distributed applications, maintaining reliability through continuous refinement. Eager to continuously expand technical knowledge and apply innovative approaches to advance engineering practices and team goals.
        \end{onecolentry}
        

    \section{Experience}
        \begin{flextwocolentry}{\textscc{August 2018 – Present}}
            \boldblack{Vermantia}, Marousi - Athens
        \end{flextwocolentry}
        \begin{onecolentry}
            \textscc{Senior Software Engineer} 
            \begin{highlights}
                \item Designed and built Vermantia’s custom digital‑signage platform, including a Django‑based CMS, a Python/FastAPI agent, and a packaging pipeline. Integrated low‑level OS components and video‑stream handling with window managers to support scalable, multi‑screen content distribution. \\ \boldblack{Key technologies}: \textit{Python}, \textit{Django}, \textit{Celery}, \textit{FastAPI}, \textit{RabbitMQ}, \textit{GObject}, \textit{GStreamer}, \textit{WebKit} and \textit{Wayland}.
                \item Transitioned deployment strategies to \textit{Kubernetes}, eliminating the need for night-time deployments while ensuring reliable, seamless releases.
                \item Designed and delivered end‑to‑end full‑stack systems for OPAP’s sports and racing channels, leveraging \textit{Django} and \textit{FastAPI} on the back end and React on the front end. Ensured continuous availability and seamless high‑traffic user experiences by integrating live streaming and real‑time data pipelines.
                \item Transitioned a \textit{.NET}‑based microservice ecosystem that ingested data updates and routed them to clients according to contractual rules into a \textit{Python} stack using \textit{Django} and \textit{Celery} pipelines, streamlining development and improving scalability.
            \end{highlights}
        \end{onecolentry}
        
        \begin{onecolentry}
            \textscc{Software Engineer} \begin{highlights}
                \item Developed and sustained an automated video‑stream orchestration platform using Django and Celery workers, eliminating the need for manual operational intervention.
                \item Optimized a Scikit‑learn‑based machine‑learning service that links sports and racing events to live streams, continuously boosting prediction accuracy.
                \item Designed and rolled out an organization‑wide logging platform, employing Loki for log aggregation and Grafana for unified visual dashboards, to provide real‑time system‑health monitoring across all teams.
                \item Re‑architected legacy microservices into a fault‑tolerant framework, markedly improving system stability and reliability.
            \end{highlights}
            
        \end{onecolentry}


    \section{Skills}
        \begin{onecolentry}
            \boldblack{Languages:} Python, TypeScript, Javascript, Golang, Java, PHP, Bash\end{onecolentry}
        \vspace{0.2 cm}      
        \begin{onecolentry}
            \boldblack{Backend Frameworks:} Django, FastAPI, Flask, Celery, SQLAlchemy, NextJS\end{onecolentry}
        \vspace{0.2 cm}
        \begin{onecolentry}
            \boldblack{Frontend Frameworks:} NextJS, React, NodeJS, jQuery\end{onecolentry}
        \vspace{0.2 cm}
        \begin{onecolentry}
            \boldblack{Devops \& Infrastructure:} Docker, Kubernetes, Terraform, Ansible, Gitlab-CI\end{onecolentry}
        \vspace{0.2 cm}
        \begin{onecolentry}
            \boldblack{Databases \& MQ:} PostgreSQL, MongoDB, RabbitMQ, Redis, Memcache\end{onecolentry}
        \vspace{0.2 cm}
        \begin{onecolentry}
            \boldblack{Testing \& Development Practices:} Git, TDD, Unittests, Integration tests\end{onecolentry}
        \vspace{0.2 cm}
        \begin{onecolentry}
            \boldblack{Operating Systems \& Virtualization :} Ubuntu, Debian, Proxmox, QEMU, HyperV\end{onecolentry}
        \vspace{0.2 cm}
        \begin{onecolentry}
            \boldblack{Networking:} Caddy, Traefik, DNS, Wireguard\end{onecolentry}
        \vspace{0.2 cm}
    
    \section{Education}
        
        \begin{threecolentry}{\boldblack{BS}}{
            \textscc{Sept. 2014 – Sept. 2019}
        }
            \boldblack{University of Piraeus}, Department of Informatics - Digital Systems
        \end{threecolentry}
        \begin{onecolentry}
                        \begin{highlights}
                \item \boldblack{Thesis:}  \textsc{Parallel execution of Reverse Top-K queries in main memory.} \href{https://gitlab.com/yagdo/thesis/raw/master/thesis.pdf?inline=false}{PDF} \\
                Conducted original research on optimizing reverse top-k queries through parallel processing in main memory, designing a parallel processing framework and programmable API to advance query efficiency in distributed systems, demonstrating scalable performance gains over single-threaded approaches.
            \end{highlights}
        \end{onecolentry}

        
    \section{Certifications}
        \begin{certtwocolentry}{
        \begin{tabular}{r@{}}
            \textscc{February 2026} \href{https://www.udemy.com/certificate/UC-c8b56189-e9b9-41c8-84d9-af9698696ea7/}{Certificate} 
            \end{tabular}
        }
            \boldblack{TypeScript:} Understanding TypeScript
        \end{certtwocolentry}        
        \vspace{0.15 cm}
        
        \begin{certtwocolentry}{
        \begin{tabular}{r@{}}
            \textscc{February 2026} \href{https://www.udemy.com/certificate/UC-67438d23-0555-4437-81b7-dc30ea0ac5d1/}{Certificate} 
            \end{tabular}
        }
            \boldblack{Ansible:} Ansible for the Absolute Beginner - Hands-On - DevOps
        \end{certtwocolentry}          
        \vspace{0.15 cm}
        
        \begin{certtwocolentry}{
        \begin{tabular}{r@{}}
             \textscc{April 2024} \href{https://www.udemy.com/certificate/UC-d9518e9d-00e4-45f9-b702-c051af042f8d/}{Certificate}
             \end{tabular}
            
        }
            \boldblack{Kubernetes:} Learn DevOps: The Complete Kubernetes Course
        \end{certtwocolentry}        
        \vspace{0.15 cm}
        
        \begin{certtwocolentry}{
        \begin{tabular}{r@{}}
            \textscc{October 2022} \href{https://www.udemy.com/certificate/UC-b7f728e8-c87f-46d5-8a06-269961fa82eb/}{Certificate} 
            \end{tabular}
        }
            \boldblack{Terraform:} AWS Infrastructure as Code With Terraform
        \end{certtwocolentry}

    \section{Projects}
        \begin{twocolentry}{
            \textscc{Ongoing} \hrefWithoutArrow{https://papaioannou.dev/}{papaioannou.dev}
        }
            \boldblack{Homelab}
            \begin{highlights}
                \item In my spare time, I've been developing a Homelab project that operates on a Proxmox cluster. This showcases my skills in designing, implementing, and managing complex systems and infrastructure. The cluster hosts various open-source services that I utilize daily.
                \item \boldblack{Technologies used:} Proxmox, Wireguard, Vaultwarden, Immich, Adguard, Open-WebUI, Ghost
            \end{highlights}
        
        \end{twocolentry}
        \vspace{0.2 cm}


        \begin{twocolentry}{
            \begin{tabular}{r@{}}
  \textscc{October 2022} \\
  \mbox{\hrefWithoutArrow{https://github.com/giannis-papaioannou/warden}{{\footnotesize\faGithub}\hspace*{0.13cm}warden}}
\end{tabular}
                
        }
            \boldblack{Warden IP tracker}
            \begin{highlights}
                \item Containerized agent that monitors dynamic IP changes and sends real-time notifications via Discord webhooks.
                \item \boldblack{Technologies used:} Python, Docker

            \end{highlights}
        \end{twocolentry}
        \vspace{0.2 cm}

        \begin{twocolentry}{
            \begin{tabular}{r@{}}
  \textscc{November 2020} \\
\mbox{\hrefWithoutArrow{https://github.com/giannis-papaioannou/steward}{{\footnotesize\faGithub}\hspace*{0.13cm}steward}}
\end{tabular}
        }
            \boldblack{Steward}
            \begin{highlights}
                \item A Docker-based service that assesses network speed and latency.
                \item \boldblack{Technologies used:} Python, Docker, InfluxDB
            \end{highlights}
        \end{twocolentry}
        \vspace{0.2 cm}

        \begin{twocolentry}{
        
            \begin{tabular}{r@{}}
  \textscc{August 2018} \\
\mbox{\hrefWithoutArrow{https://github.com/giannis-papaioannou/scarletmoon}{{\footnotesize\faGithub}\hspace*{0.13cm}scarletmoon}}
\end{tabular}
        }
            \boldblack{Scarletmoon}
            \begin{highlights}
                \item A service that creates inverted indexes to facilitate efficient text searching without needing predefined document structures.
                \item \boldblack{Technologies used:} Python, Nameko
            \end{highlights}
        \end{twocolentry}
        \vspace{0.2 cm}
        
        \begin{twocolentry}{
        
        
            \begin{tabular}{r@{}}
  \textscc{February 2018} \\
  \mbox{\hrefWithoutArrow{https://github.com/giannis-papaioannou/allpayauction}{{\footnotesize\faGithub}\hspace*{0.13cm}allpayauction}}%
\end{tabular}
        }
            \boldblack{All Pay Auction}
            \begin{highlights}
                \item This is a university project for the algorithmic trading course. It involves implementing an iterative best response method, determining the optimal move by taking into account the actions of other bidders.
                \item \boldblack{Technologies used:} Python, Matplotlib
            \end{highlights}
        \end{twocolentry}
        \vspace{0.2 cm}
        
        \begin{twocolentry}{
        
        
            \begin{tabular}{r@{}}
  \textscc{July 2017} \\
  
                \mbox{\hrefWithoutArrow{https://github.com/giannis-papaioannou/flamel}{{\footnotesize\faGithub}\hspace*{0.13cm}flamel}}%
\end{tabular}
        }
            \boldblack{Flamel}
            \begin{highlights}
                \item Graphical interface showcasing different path-finding algorithms.
                \item \boldblack{Technologies used:} Java
            \end{highlights}
        \end{twocolentry}
        \vspace{0.2 cm}
    

\end{document}